\section{Machine Learning and Deep Learning}
\label{sec:ml_intro}

Machine Learning (ML), and in particular \emph{Deep Learning} (DL), is the field of study concerned with designing and training parameterized mathematical models to perform tasks through experience rather than explicit programming \cite{heaton_ian_2018,noauthor_pattern_2007}. Learning algorithms can be categorized in various ways, two of the most prominent paradigms being \emph{supervised} and \emph{unsupervised} learning. 

In general, a machine learning model is a parameterized function 
\begin{equation}
    f_{\theta} : \mathcal{X} \rightarrow \mathcal{Y},
\end{equation}
where $\mathcal{X}$ denotes the input domain (often referred to as the \emph{feature space}), $\mathcal{Y}$ the output domain (labels or targets), and $\theta \in \Theta$ the set of \emph{trainable parameters}. The model is evaluated according to a task-specific \emph{loss function} 
\begin{equation}
    \mathcal{L} : \Theta \times \mathcal{X} \rightarrow \mathbb{R},
\end{equation}
which measures the ability of the model with parameters $\theta$ to perform the desired task on an input $x \in \mathcal{X}$. For convenience, we assume that smaller values of $\mathcal{L}$ correspond to better performance. The goal of the training process is to find parameters $\theta^*$ that minimize this loss, typically formulated as the optimization problem
\begin{equation}
    \label{equ:training_task}
    \theta^* = \argmin_{\theta \in \Theta} \, \mathbb{E}_{x \sim \mathcal{D}} \, \mathcal{L}(\theta, x),
\end{equation}
where $\mathcal{D}$ denotes the data distribution over the input space.

In \emph{supervised learning}, the dataset $\mathcal{D}$ provides labeled examples $(x, y)$, allowing direct estimation of $\theta^*$ by gradient-based optimization (e.g., stochastic gradient descent). In contrast, \emph{unsupervised learning} seeks to discover hidden structure in unlabeled data, often by optimizing surrogate objectives such as likelihoods or reconstruction errors. Many modern frameworks extend these two paradigms to encompass reinforcement learning, self-supervised learning, and generative modeling, broadening the scope of ML beyond prediction toward reasoning and structured decision-making \cite{lecun_deep_2015}.

In this work, we focus on learning paradigms where ML models are used to approximate or guide optimization processes, often by learning mappings from structured problem instances to high-quality solutions or to efficient inference strategies. Such approaches are particularly relevant in the context of combinatorial optimization and probabilistic graphical models, as discussed in later sections (see \S\ref{sec:mrf_intro} and \S\ref{sec:co_intro}).


\subsection{Supervised Learning}
\label{sec:supervised_learning}

Supervised learning refers to the class of learning strategies in which the \emph{ground truth} or target value is available during training. Formally, we assume access to a dataset of input–output pairs 
\begin{equation}
    \mathcal{D} = \{(x_i, y_i)\}_{i=1}^{N}, \quad x_i \in \mathcal{X}, \; y_i \in \mathcal{Y},
\end{equation}
where each input $x_i$ is associated with a corresponding label $y_i$ representing the correct or desired outcome. The objective is to learn a parameterized function $f_{\theta} : \mathcal{X} \rightarrow \mathcal{Y}$ capable of producing predictions $\hat{y} = f_{\theta}(x)$ that generalize well to unseen data drawn from the same underlying distribution $\mathcal{D}$.

By quantifying the discrepancy between the predicted outcome $\hat{y}$ and the ground truth $y$, it becomes possible to define a \emph{loss function} $\mathcal{L}_{\text{sup}} : \Theta \times \mathcal{X} \times \mathcal{Y} \rightarrow \mathbb{R}$ that measures model performance on each training example. The learning problem is then formulated as minimizing the expected loss over the data distribution:
\begin{equation}
    \label{equ:supervised_learning}
    \theta^* = \argmin_{\theta \in \Theta} \, \mathbb{E}_{(x,y) \sim \mathcal{D}} \left[ \mathcal{L}_{\text{sup}}(\theta, x, y) \right].
\end{equation}

In practice, this expectation is approximated by the empirical average over the training set:
\begin{equation}
    \widehat{\mathcal{L}}_{\text{sup}}(\theta) = \frac{1}{N} \sum_{i=1}^{N} \mathcal{L}_{\text{sup}}(\theta, x_i, y_i).
\end{equation}

The form of $\mathcal{L}_{\text{sup}}$ depends on the prediction task:
\begin{itemize}
    \item In \emph{regression} tasks, the model predicts continuous quantities and typical losses include the squared error $\mathcal{L}_{\text{sup}} = \| f_{\theta}(x) - y \|^2$ or the mean absolute error.
    \item In \emph{classification} tasks, where the label $y$ belongs to a discrete set of classes, the loss commonly takes the form of the negative log-likelihood or cross-entropy:
    \begin{equation}
        \mathcal{L}_{\text{sup}}(\theta, x, y) = - \log p_{\theta}(y|x),
    \end{equation}
    where $p_{\theta}(y|x)$ denotes the model’s predicted probability for the correct class.
\end{itemize}

Supervised learning thus reduces to a parameter optimization problem, where $\theta$ is iteratively updated to minimize $\widehat{\mathcal{L}}_{\text{sup}}(\theta)$ using stochastic gradient descent or one of its many variants \cite{goodfellow_deep_2016,kingma_adam_2014}. When the learned parameters generalize beyond the training data, the model effectively captures an underlying mapping between inputs and outputs, a property central to most modern applications of deep learning \cite{vapnik_nature_2000,bishop_pattern_2006}.

\subsection{Unsupervised Learning}
\label{sec:unsupervised_learning}

Unsupervised learning encompasses learning paradigms in which no explicit ground-truth labels or target outputs are provided during training. Instead, the objective is to discover latent structures, dependencies, or representations within the data itself \cite{bishop_pattern_2006,goodfellow_deep_2016}. Formally, given a dataset of unlabeled examples 
\begin{equation}
    \mathcal{D} = \{ x_i \}_{i=1}^{N}, \quad x_i \in \mathcal{X},
\end{equation}
the goal is to learn parameters $\theta$ that optimize a self-consistency criterion defined by a loss function $\mathcal{L}_{\text{unsup}} : \Theta \times \mathcal{X} \rightarrow \mathbb{R}$,
\begin{equation}
    \label{equ:unsupervised_learning}
    \theta^* = \argmin_{\theta \in \Theta} \, \mathbb{E}_{x \sim \mathcal{D}} \left[ \mathcal{L}_{\text{unsup}}(\theta, x) \right].
\end{equation}

The precise form of $\mathcal{L}_{\text{unsup}}$ depends on the learning objective. Common unsupervised frameworks include:

\paragraph{1. Representation learning and dimensionality reduction.}
Methods such as \emph{Principal Component Analysis (PCA)} \cite{jolliffe_principal_2010}, \emph{autoencoders} \cite{hinton_reducing_2006}, and their probabilistic extensions aim to learn low-dimensional latent representations $z = g_{\theta}(x)$ that capture the salient structure of the data. In an autoencoder, for instance, the model is composed of an encoder–decoder pair $(g_{\theta}, d_{\theta})$ trained to minimize the reconstruction error:
\begin{equation}
    \mathcal{L}_{\text{rec}}(\theta, x) = \| d_{\theta}(g_{\theta}(x)) - x \|^2,
\end{equation}
thus encouraging the latent code $z$ to capture the essential variability in $\mathcal{X}$ while discarding noise and redundancy \cite{ hinton_reducing_2006,kingma_auto-encoding_2022}.

\paragraph{2. Density estimation and probabilistic modeling.}
Another key class of unsupervised methods seeks to model the underlying data distribution $p_{\theta}(x)$ directly. Training such models typically involves maximizing the likelihood of the observed data, or equivalently, minimizing the negative log-likelihood loss:
\begin{equation}
    \mathcal{L}_{\text{nll}}(\theta, x) = - \log p_{\theta}(x).
\end{equation}
Examples include probabilistic graphical models such as Markov Random Fields (MRFs) and Conditional Random Fields (CRFs), as well as modern generative models like Variational Autoencoders (VAEs) \cite{kingma_introduction_2019} and Generative Adversarial Networks (GANs) \cite{goodfellow_generative_2014}. These models can capture complex dependencies between variables and generate realistic new samples from learned distributions.

\paragraph{3. Clustering and self-organization.}
In clustering-based approaches, the learning process groups similar samples together according to an implicit or explicit similarity measure. Classical
algorithms include $k$-means, Gaussian mixture models, and spectral clustering \cite{macqueen_methods_1967,dempster_maximum_1977,ng_spectral_2001}. In the neural setting, differentiable clustering losses or contrastive objectives are often employed to encourage representation spaces where semantically similar samples are close together \cite{oord_representation_2019,chen_simple_2020}.

Overall, unsupervised learning provides a foundation for learning useful data representations and priors that can later be exploited in semi-supervised, transfer, or reinforcement learning contexts. In this work, it plays a key role in formulating models that learn from problem structure—rather than explicit supervision—especially in the context of probabilistic graphical models and combinatorial optimization, where ground-truth solutions are expensive or unavailable.


\subsection{Generative Models and Denoising Objectives}

Generative models aim to learn a probability distribution over the data, rather than directly predicting a single deterministic output. In the conditional setting considered throughout this thesis, an input instance $x \in \mathcal{X}$ may be associated with multiple valid or high-quality outputs $y \in \mathcal{Y}$. This is particularly relevant in combinatorial optimization, where a problem instance can admit several near-optimal solutions with different structures. Instead of returning a single prediction, a generative model defines a conditional distribution over possible outputs and obtains predictions by sampling from this distribution (see Figure~\ref{fig:conditional_generative_graph}).

Formally, we view input-output pairs as realizations of random variables $X$ and $Y$. The goal is to approximate the conditional distribution
\begin{equation}
    p(y \mid x) = \mathbb{P}\left[ Y = y \mid X = x \right],
\end{equation}
where $x$ denotes the observed input instance and $y$ denotes a target output, such as a feasible solution or an optimal assignment.

In practice, the true conditional distribution $p(y \mid x)$ is unknown and must be approximated by a parameterized model. We denote this approximation by $p_\Theta(y \mid x)$, where $\Theta$ are learned parameters. In some models, the input $x$ is first mapped to the parameters of an output distribution, for example through a neural network $\Theta = f_\theta(x)$, yielding a distribution $p_\Theta(y)$. Predictions can then be obtained by sampling
\begin{equation}
    y \sim p_\Theta(\cdot \mid x),
\end{equation}
or by selecting high-probability outputs according to the learned distribution.

Different families of generative models mainly differ in how they parameterize this distribution and in the training objective used to match it to the data distribution. Examples include maximum-likelihood objectives, variational objectives, adversarial objectives, denoising objectives, and, more recently, diffusion and flow-matching objectives. In the remainder of this section, we focus on denoising-based generative models, since they provide the conceptual basis for the diffusion and flow-based decoding methods studied later in this thesis.

\begin{figure}[t]
\centering
\begin{tikzpicture}[
    font=\small,
    >=Latex,
    panel/.style={
        rounded corners=8pt,
        draw=#1,
        fill=#1!6,
        line width=0.9pt,
        inner sep=6pt
    },
    stagecircle/.style={
        circle,
        fill=#1,
        text=white,
        font=\bfseries,
        minimum size=7mm,
        inner sep=0pt
    },
    arrow/.style={
        -{Latex[length=3mm,width=2mm]},
        line width=1.1pt,
        draw=blue!55!black
    },
    dottedarrow/.style={
        -{Latex[length=3mm,width=2mm]},
        line width=1.1pt,
        draw=blue!55!black,
        dotted
    },
    itemrow/.style={
        rounded corners=3pt,
        draw=gray!30,
        fill=white,
        minimum height=6mm,
        inner sep=2pt
    },
    probbar/.style={
        draw=none,
        fill=green!60!black,
        rounded corners=1pt
    },
    nodecircle/.style={
        circle,
        draw=blue!50!black,
        fill=blue!15,
        minimum size=3.5mm,
        inner sep=0pt
    },
    hiddennode/.style={
        circle,
        draw=green!50!black,
        fill=green!25,
        minimum size=3.5mm,
        inner sep=0pt
    },
    outputnode/.style={
        circle,
        draw=purple!60!black,
        fill=purple!20,
        minimum size=3.5mm,
        inner sep=0pt
    }
]

% =========================================================
% Colors
% =========================================================
\definecolor{TitleColor}{RGB}{10,25,70}
\definecolor{PurplePanel}{RGB}{130,100,200}
\definecolor{BluePanel}{RGB}{65,130,200}
\definecolor{GreenPanel}{RGB}{40,135,85}
\definecolor{OrangePanel}{RGB}{220,145,25}

\definecolor{NodeA}{RGB}{220,50,47}
\definecolor{NodeB}{RGB}{38,139,210}
\definecolor{NodeC}{RGB}{133,153,0}
\definecolor{NodeD}{RGB}{181,137,0}
\definecolor{NodeE}{RGB}{108,113,196}

% =========================================================
% Panel coordinates
% =========================================================
\coordinate (P1) at (-6.0,0);
\coordinate (P2) at (-2.0,0);
\coordinate (P3) at ( 2.1,0);
\coordinate (P4) at ( 6.1,0);

% =========================================================
% Stage badges and headings (moved higher)
% =========================================================
% \node[stagecircle=PurplePanel] at ($(P1)+(0,2.65)$) {1};
% \node[font=\bfseries\large, text=TitleColor] at ($(P1)+(0,2.25)$) {Input};

% \node[stagecircle=BluePanel] at ($(P2)+(0,2.65)$) {2};
% \node[font=\bfseries\large, text=TitleColor] at ($(P2)+(0,2.25)$) {Model};

% \node[stagecircle=GreenPanel] at ($(P3)+(0,2.65)$) {3};
% \node[font=\bfseries\large, text=TitleColor] at ($(P3)+(0,2.25)$) {Distribution};

% \node[stagecircle=OrangePanel] at ($(P4)+(0,2.65)$) {4};
% \node[font=\bfseries\large, text=TitleColor] at ($(P4)+(0,2.25)$) {Sample};

% =========================================================
% Stage 1: Input graph
% =========================================================
\node[panel=PurplePanel, minimum width=3.0cm, minimum height=4.3cm, anchor=center] (inputpanel) at (P1) {};
\node[font=\bfseries, text=TitleColor] at ($(P1)+(0,1.75)$) {Input};

\begin{scope}[shift={(P1)}, scale=0.92]
    % Graph nodes
    \node[circle, draw=black, fill=NodeA!75, minimum size=7mm, inner sep=0pt] (gA) at (-0.55,0.95) {};
    \node[circle, draw=black, fill=NodeB!75, minimum size=7mm, inner sep=0pt] (gB) at (0.55,1.05) {};
    \node[circle, draw=black, fill=NodeC!75, minimum size=7mm, inner sep=0pt] (gC) at (-0.95,-0.05) {};
    \node[circle, draw=black, fill=NodeD!75, minimum size=7mm, inner sep=0pt] (gD) at (0.05,-0.15) {};
    \node[circle, draw=black, fill=NodeE!75, minimum size=7mm, inner sep=0pt] (gE) at (0.95,-0.85) {};

    % Graph edges
    \draw[thick, gray!70] (gA) -- (gB);
    \draw[thick, gray!70] (gA) -- (gC);
    \draw[thick, gray!70] (gA) -- (gD);
    \draw[thick, gray!70] (gB) -- (gD);
    \draw[thick, gray!70] (gC) -- (gD);
    \draw[thick, gray!70] (gD) -- (gE);

    % Labels
    \node[font=\scriptsize] at (gA) {$1$};
    \node[font=\scriptsize] at (gB) {$2$};
    \node[font=\scriptsize] at (gC) {$3$};
    \node[font=\scriptsize] at (gD) {$4$};
    \node[font=\scriptsize] at (gE) {$5$};
\end{scope}

\node[font=\Large] at ($(P1)+(0,-1.5)$) {$x$};

% =========================================================
% Stage 2: Model panel
% =========================================================
\node[panel=BluePanel, minimum width=3.5cm, minimum height=4.3cm] (modelpanel) at (P2) {};
\node[font=\bfseries, text=TitleColor] at ($(P2)+(0,1.75)$) {Neural network};

\begin{scope}[shift={(P2)}, scale=0.88]
    % Neural network nodes
    \foreach \i/\yy in {1/0.78,2/0.33,3/-0.12,4/-0.57}{
        \node[nodecircle] (in\i) at (-0.95,\yy) {};
    }
    \foreach \i/\yy in {1/0.82,2/0.37,3/-0.08,4/-0.53}{
        \node[hiddennode] (hA\i) at (-0.25,\yy) {};
        \node[hiddennode] (hB\i) at (0.42,\yy) {};
    }
    \foreach \i/\yy in {1/0.58,2/0.02,3/-0.54}{
        \node[outputnode] (out\i) at (1.02,\yy) {};
    }

    \foreach \i in {1,2,3,4}{
        \foreach \j in {1,2,3,4}{
            \draw[blue!45!black, line width=0.35pt] (in\i) -- (hA\j);
            \draw[blue!45!black, line width=0.35pt] (hA\i) -- (hB\j);
        }
    }
    \foreach \i in {1,2,3,4}{
        \foreach \j in {1,2,3}{
            \draw[blue!45!black, line width=0.35pt] (hB\i) -- (out\j);
        }
    }
\end{scope}

\node[font=\Large] at ($(P2)+(0,-1.45)$) {$p_\theta(y \mid x)$};

% =========================================================
% Stage 3: Distribution panel
% =========================================================
\node[panel=GreenPanel, minimum width=3.25cm, minimum height=4.3cm] (distpanel) at (P3) {};
\node[font=\bfseries, text=TitleColor, align=center] at ($(P3)+(0,1.48)$) {Node selection\\probabilities};

\begin{scope}[shift={(P3)}, scale=0.88]
    \foreach \y/\name/\col/\prob/\barlen in {
        0.82/Node 1/NodeA/0.82/1.05,
        0.35/Node 2/NodeB/0.21/0.27,
        -0.12/Node 3/NodeC/0.73/0.92,
        -0.59/Node 4/NodeD/0.15/0.20,
        -1.06/Node 5/NodeE/0.78/0.98
    }{
        \node[itemrow, minimum width=2.65cm, minimum height=3.8mm] at (0,\y) {};
        \node[circle, draw=black, fill=\col!75, minimum size=2.8mm, inner sep=0pt] at (-1.02,\y) {};
        %\node[anchor=west, font=\scriptsize] at (-0.82,\y) {\name};
        \draw[probbar] (0.08,\y-0.08) rectangle ++(\barlen,0.16);
        %\node[anchor=east, font=\scriptsize, text=GreenPanel!80!black] at (1.25,\y) {\prob};
    }
\end{scope}

%\node[font=\large] at ($(P3)+(0,-1.55)$) {$y \sim p_\theta(\cdot \mid x)$};

% =========================================================
% Stage 4: Sampled solution graph
% =========================================================
\node[panel=OrangePanel, minimum width=3.0cm, minimum height=4.3cm] (samplepanel) at (P4) {};
\node[font=\bfseries, text=TitleColor] at ($(P4)+(0,1.75)$) {Sample};

\begin{scope}[shift={(P4)}, scale=0.92]
    % Faded original graph
    \node[circle, draw=gray!60, fill=NodeA!20, minimum size=7mm, inner sep=0pt] (sA) at (-0.55,0.95) {};
    \node[circle, draw=gray!60, fill=gray!15, minimum size=7mm, inner sep=0pt] (sB) at (0.55,1.05) {};
    \node[circle, draw=gray!60, fill=NodeC!25, minimum size=7mm, inner sep=0pt] (sC) at (-0.95,-0.05) {};
    \node[circle, draw=gray!60, fill=gray!15, minimum size=7mm, inner sep=0pt] (sD) at (0.05,-0.15) {};
    \node[circle, draw=gray!60, fill=NodeE!25, minimum size=7mm, inner sep=0pt] (sE) at (0.95,-0.85) {};

    \draw[thick, gray!35] (sA) -- (sB);
    \draw[thick, gray!35] (sA) -- (sC);
    \draw[thick, gray!35] (sA) -- (sD);
    \draw[thick, gray!35] (sB) -- (sD);
    \draw[thick, gray!35] (sC) -- (sD);
    \draw[thick, gray!35] (sD) -- (sE);

    % Highlight sampled nodes only
    \draw[line width=1.2pt, black] (sA) circle (3.6mm);
    \draw[line width=1.2pt, black] (sC) circle (3.6mm);
    \draw[line width=1.2pt, black] (sE) circle (3.6mm);

    \node[font=\scriptsize] at (sA) {$1$};
    \node[font=\scriptsize, text=gray!60] at (sB) {$2$};
    \node[font=\scriptsize] at (sC) {$3$};
    \node[font=\scriptsize, text=gray!60] at (sD) {$4$};
    \node[font=\scriptsize] at (sE) {$5$};
\end{scope}

\node[font=\large] at ($(P4)+(0,-1.5)$)
{$\{1,3,5\}$};

% =========================================================
% Arrows between stages
% =========================================================
\draw[arrow] ($(inputpanel.east)+(0.10,0)$) -- ($(modelpanel.west)+(-0.10,0)$);
\draw[arrow] ($(modelpanel.east)+(0.10,0)$) -- ($(distpanel.west)+(-0.10,0)$);
\draw[dottedarrow] ($(distpanel.east)+(0.10,0)$) -- ($(samplepanel.west)+(-0.10,0)$)
    node[midway, above, font=\scriptsize\itshape, text=blue!55!black] {};

\end{tikzpicture}
\caption{Illustration of conditional generative modeling for a combinatorial optimization problem. A graph-structured input instance is processed by a neural network, which predicts selection probabilities for the nodes. A candidate solution is then obtained by sampling a subset of nodes from the predicted distribution.}
\label{fig:conditional_generative_graph}
\end{figure}

\paragraph{Denoising-based generative modeling.}

A central idea in modern generative modeling is to learn how to recover structured data from corrupted observations. Instead of directly learning to sample from the data distribution, denoising-based methods define a corruption process that progressively perturbs a clean target $y$ into a noisy version $\tilde{y}$. A model is then trained to invert this process, either by predicting the original clean sample, the noise that was added, or an update direction that moves the corrupted sample closer to the data distribution.

Let $y \in \mathcal{Y}$ denote a clean target output, such as an optimal or high-quality solution of a combinatorial optimization problem. A corruption process defines a conditional distribution
\begin{equation}
    q(\tilde{y} \mid y, t),
\end{equation}
where $t$ denotes the noise level or time index. For small values of $t$, the corrupted sample $\tilde{y}$ remains close to the original target, while for large values of $t$ it may contain little information about $y$. The learning problem consists in training a model that, given the corrupted sample $\tilde{y}$, the input instance $x$, and the noise level $t$, predicts how to recover information about the clean output:
\begin{equation}
    D_\theta(x, \tilde{y}, t) \approx y.
\end{equation}
Equivalently, depending on the formulation, the model may predict the perturbation applied to $y$, or a direction in the output space that transforms $\tilde{y}$ toward a clean sample.

\begin{figure}[t]
\centering
\begin{tikzpicture}[
    font=\small,
    >=Latex,
    arrow/.style={
        -{Latex[length=3mm,width=2mm]},
        line width=1.1pt,
        draw=blue!60!black
    },
    panel/.style={
        rounded corners=8pt,
        line width=0.9pt,
        inner sep=6pt
    },
    stagebadge/.style={
        circle,
        minimum size=7mm,
        inner sep=0pt,
        font=\bfseries,
        text=white
    },
    clean/.style={
        circle,
        draw=black,
        line width=0.8pt,
        minimum size=7mm,
        inner sep=0pt
    },
    noisy/.style={
        circle,
        draw=gray!70,
        fill=gray!30,
        line width=0.8pt,
        minimum size=7mm,
        inner sep=0pt
    },
    inactive/.style={
        circle,
        draw=gray!60,
        fill=white,
        line width=0.8pt,
        minimum size=7mm,
        inner sep=0pt
    },
    nninput/.style={
        circle,
        draw=blue!55!black,
        fill=blue!15,
        minimum size=3.8mm,
        inner sep=0pt
    },
    nnhidden/.style={
        circle,
        draw=green!50!black,
        fill=green!25,
        minimum size=3.8mm,
        inner sep=0pt
    },
    nnoutput/.style={
        circle,
        draw=purple!60!black,
        fill=purple!20,
        minimum size=3.8mm,
        inner sep=0pt
    }
]

% -------------------------------------------------
% Colors
% -------------------------------------------------
\definecolor{TitleColor}{RGB}{10,25,70}
\definecolor{PurplePanel}{RGB}{130,100,200}
\definecolor{BluePanel}{RGB}{65,130,200}
\definecolor{GreenPanel}{RGB}{40,135,85}

\definecolor{NodeSel}{RGB}{166,206,57}   % selected node
\definecolor{NodeAlt}{RGB}{102,194,255}  % alternate highlight
\definecolor{NodeWarn}{RGB}{255,180,70}  % optional accent

% -------------------------------------------------
% Coordinates of the three stages
% -------------------------------------------------
\coordinate (P1) at (-4.8,0);
\coordinate (P2) at (0,0);
\coordinate (P3) at (4.8,0);

% -------------------------------------------------
% Headings
% -------------------------------------------------
% \node[stagebadge, fill=PurplePanel] at ($(P1)+(0,2.5)$) {1};
% \node[font=\bfseries\large, text=TitleColor] at ($(P1)+(0,2.1)$) {Noisy input};

% \node[stagebadge, fill=BluePanel] at ($(P2)+(0,2.5)$) {2};
% \node[font=\bfseries\large, text=TitleColor] at ($(P2)+(0,2.1)$) {Prediction};

% \node[stagebadge, fill=GreenPanel] at ($(P3)+(0,2.5)$) {3};
% \node[font=\bfseries\large, text=TitleColor] at ($(P3)+(0,2.1)$) {Denoised output};

% -------------------------------------------------
% Panel 1 : Noisy solution
% -------------------------------------------------
\node[panel, draw=PurplePanel, fill=PurplePanel!6, minimum width=3.3cm, minimum height=4.1cm] (noisypanel) at (P1) {};

\node[font=\bfseries, text=TitleColor] at ($(P1)+(0,1.75)$) {Noise};

\begin{scope}[shift={(P1)}, scale=0.95]
    % nodes
    \node[clean, fill=NodeSel!85] (a1) at (-0.75, 0.95) {};
    \node[noisy]               (a2) at ( 0.35, 1.05) {};
    \node[clean, fill=NodeSel!85] (a3) at (-1.05,-0.05) {};
    \node[noisy]               (a4) at ( 0.10,-0.18) {};
    \node[inactive]            (a5) at ( 0.95,-0.75) {};
    \node[inactive]            (a6) at (-0.20,-0.95) {};

    % edges
    \draw[thick, gray!65] (a1) -- (a2);
    \draw[thick, gray!65] (a1) -- (a3);
    \draw[thick, gray!65] (a1) -- (a4);
    \draw[thick, gray!65] (a2) -- (a4);
    \draw[thick, gray!65] (a3) -- (a4);
    \draw[thick, gray!65] (a4) -- (a5);
    \draw[thick, gray!65] (a3) -- (a6);

    % labels
    \node[font=\scriptsize] at (a1) {$1$};
    \node[font=\scriptsize, text=gray!60!black] at (a2) {?};
    \node[font=\scriptsize] at (a3) {$3$};
    \node[font=\scriptsize, text=gray!60!black] at (a4) {?};
    \node[font=\scriptsize, text=gray!70!black] at (a5) {$5$};
    \node[font=\scriptsize, text=gray!70!black] at (a6) {$6$};
\end{scope}

\node[align=center] at ($(P1)+(0,-1.8)$)
{Noisy solution $\tilde y$};

% -------------------------------------------------
% Panel 2 : Denoising model
% -------------------------------------------------
\node[panel, draw=BluePanel, fill=BluePanel!6, minimum width=3.7cm, minimum height=4.1cm] (modelpanel) at (P2) {};
\node[font=\bfseries, text=TitleColor] at ($(P2)+(0,1.75)$) {Denoising model};

\begin{scope}[shift={(P2)}, scale=0.9]
    % network
    \foreach \i/\yy in {1/0.75,2/0.25,3/-0.25,4/-0.75}{
        \node[nninput] (i\i) at (-0.95,\yy) {};
    }
    \foreach \i/\yy in {1/0.9,2/0.3,3/-0.3,4/-0.9}{
        \node[nnhidden] (h\i) at (0,\yy) {};
    }
    \foreach \i/\yy in {1/0.5,2/0.0,3/-0.5}{
        \node[nnoutput] (o\i) at (0.95,\yy) {};
    }

    \foreach \u in {1,2,3,4}{
        \foreach \v in {1,2,3,4}{
            \draw[blue!50!black, line width=0.35pt] (i\u) -- (h\v);
        }
    }
    \foreach \u in {1,2,3,4}{
        \foreach \v in {1,2,3}{
            \draw[blue!50!black, line width=0.35pt] (h\u) -- (o\v);
        }
    }
\end{scope}

\node[align=center] at ($(P2)+(0,-1.8)$)
{$\hat y = f_\theta(\tilde y, x)$};

% -------------------------------------------------
% Panel 3 : Denoised output
% -------------------------------------------------
\node[panel, draw=GreenPanel, fill=GreenPanel!6, minimum width=3.3cm, minimum height=4.1cm] (cleanpanel) at (P3) {};
\node[font=\bfseries, text=TitleColor] at ($(P3)+(0,1.75)$) {Denoised};

\begin{scope}[shift={(P3)}, scale=0.95]
    % nodes
    \node[clean, fill=NodeSel!85] (b1) at (-0.75, 0.95) {};
    \node[inactive]               (b2) at ( 0.35, 1.05) {};
    \node[clean, fill=NodeSel!85] (b3) at (-1.05,-0.05) {};
    \node[inactive]               (b4) at ( 0.10,-0.18) {};
    \node[inactive]               (b5) at ( 0.95,-0.75) {};
    \node[clean, fill=NodeSel!85] (b6) at (-0.20,-0.95) {};

    % edges
    \draw[thick, gray!65] (b1) -- (b2);
    \draw[thick, gray!65] (b1) -- (b3);
    \draw[thick, gray!65] (b1) -- (b4);
    \draw[thick, gray!65] (b2) -- (b4);
    \draw[thick, gray!65] (b3) -- (b4);
    \draw[thick, gray!65] (b4) -- (b5);
    \draw[thick, gray!65] (b3) -- (b6);

    % labels
    \node[font=\scriptsize] at (b1) {$1$};
    \node[font=\scriptsize, text=gray!70!black] at (b2) {$2$};
    \node[font=\scriptsize] at (b3) {$3$};
    \node[font=\scriptsize, text=gray!70!black] at (b4) {$4$};
    \node[font=\scriptsize, text=gray!70!black] at (b5) {$5$};
    \node[font=\scriptsize] at (b6) {$6$};
\end{scope}

\node[align=center] at ($(P3)+(0,-1.8)$)
{Solution $\hat y$};

% -------------------------------------------------
% Arrows
% -------------------------------------------------
\draw[arrow] ($(noisypanel.east)+(0.08,0)$) -- ($(modelpanel.west)+(-0.08,0)$);
\draw[arrow] ($(modelpanel.east)+(0.08,0)$) -- ($(cleanpanel.west)+(-0.08,0)$);


\end{tikzpicture}
\caption{Illustration of denoising for graph-structured solutions. Starting from a corrupted or partially masked solution $\tilde y$, a prediction model infers corrected node assignments and outputs a denoised solution $\hat y$. Gray nodes indicate noisy or uncertain values.}
\label{fig:graph_denoising_illustration}
\end{figure}

This denoising view is especially natural for structured prediction. In many problems, the output is not a single scalar label but a structured object composed of many dependent decisions. In combinatorial optimization, for example, a solution may be represented by a binary vector $y \in \{0,1\}^n$, where each component corresponds to the value of a decision variable. Corrupting such a solution can be interpreted as randomly masking, flipping, or relaxing some of its variables. The denoising model must then learn how local and global dependencies between variables constrain the reconstruction of a valid or high-quality solution.

A common training objective is to minimize a reconstruction loss between the model prediction and the clean target. For example, when the model predicts the original solution, one may use
\begin{equation}
    \mathcal{L}_{\mathrm{denoise}}(\theta)
    =
    \mathbb{E}_{(x,y),t,\tilde{y}}
    \left[
        \ell\left(D_\theta(x,\tilde{y},t), y\right)
    \right],
\end{equation}
where $\tilde{y} \sim q(\cdot \mid y,t)$ and $\ell$ is a suitable loss function, such as a cross-entropy loss for discrete variables or a squared error loss for continuous relaxations. The model is therefore trained on many reconstruction tasks of varying difficulty, ranging from lightly corrupted samples to highly noisy ones.

At inference time, denoising-based generative models reverse the corruption process. Starting from a noisy or partially initialized sample, the learned denoising model is applied iteratively to produce increasingly structured outputs. This procedure can be interpreted as a learned refinement process: each denoising step uses the instance information $x$, the current candidate solution, and the noise level to move the sample toward the learned solution distribution.

In the context of this thesis, this perspective is useful because it connects generative modeling with solution construction. Rather than predicting each variable independently in a single forward pass, a denoising model defines a dynamics over candidate solutions. This makes it possible to progressively refine a noisy assignment into a feasible or high-quality solution, while potentially capturing dependencies between variables through the architecture used to parameterize the denoising function.

This general principle underlies both diffusion models, which learn to reverse a stochastic noising process, and flow-based formulations, which learn a deterministic or approximately deterministic transformation from a simple initial distribution toward the data distribution.

\paragraph{Flow matching and velocity prediction models.}

Flow matching is a generative modeling framework in which the model learns a time-dependent transformation between a simple source distribution and the data distribution. Instead of learning a reverse stochastic denoising chain, as in diffusion models, flow matching trains a neural network to predict a velocity field that transports samples from an initial distribution toward the target distribution.

Let $y_0 \sim p_0$ denote a sample from a simple source distribution, such as Gaussian noise, and let $y_1 \sim p_{\mathrm{data}}$ denote a target data sample. Flow matching defines a family of intermediate states $y_t$ for $t \in [0,1]$ that interpolate between $y_0$ and $y_1$. A common choice is the linear interpolation path
\begin{equation}
    y_t = (1-t)y_0 + t y_1.
\end{equation}
Along this path, the corresponding target velocity is
\begin{equation}
    u_t = \frac{d y_t}{dt} = y_1 - y_0.
\end{equation}
The learning problem is then to train a parameterized model $v_\theta$ to predict this velocity from the current state and time:
\begin{equation}
    v_\theta(y_t,t) \approx u_t.
\end{equation}
In a conditional setting, where generation depends on an input instance $x$, the model is instead written as
\begin{equation}
    v_\theta(x_t,t,x) \approx u_t.
\end{equation}
For example, in combinatorial optimization, $x$ may represent the optimization instance and $y_1$ may represent an optimal or high-quality solution.

The standard flow matching objective takes the form
\begin{equation}
    \mathcal{L}_{\mathrm{FM}}(\theta)
    =
    \mathbb{E}_{t,y_0,y_1}
    \left[
        \left\|
            v_\theta(y_t,t,x) - u_t
        \right\|_2^2
    \right],
\end{equation}
where $t$ is sampled from $[0,1]$, $y_0$ is sampled from the source distribution, $y_1$ is sampled from the data distribution, and $y_t$ is obtained from the chosen interpolation path. The model is therefore trained by supervised regression onto a known velocity target.

Once the velocity field has been learned, samples can be generated by starting from $y_0 \sim p_0$ and solving the ordinary differential equation (see Algorithm~\ref{alg:flow_sampling})
\begin{equation}
    \frac{d y_t}{dt} = v_\theta(y_t,t,x),
    \qquad y_0 \sim p_0.
\end{equation}
The endpoint $y_1$ of this trajectory is then expected to follow the learned data distribution. In practice, this ODE is solved numerically using a finite number of integration steps. Compared with diffusion models, this formulation often leads to a more direct generative procedure, since the model predicts a deterministic direction of motion rather than the parameters of a stochastic reverse transition.

\begin{algorithm}[t]
\caption{Sample generation with a learned velocity field}
\label{alg:flow_sampling}
\DontPrintSemicolon

\KwIn{
    Problem instance $x$;
    learned velocity field $v_\theta$;
    source distribution $p_0$;
    number of integration steps $K$;
    decoding operator $\mathcal{D}$
}
\KwOut{
    Discrete candidate solution $\hat{y}$
}

Sample an initial state $y^{(0)} \sim p_0$\;
Set $\Delta t \leftarrow 1/K$\;

\For{$k = 0,\ldots,K-1$}{
    Set $t_k \leftarrow k \Delta t$\;
    
    Predict the velocity $\hat{u}^{(k)} \leftarrow v_\theta(y^{(k)}, t_k, x)$\;
    
    Update state using an Euler step $y^{(k+1)} \leftarrow y^{(k)} + \Delta t \, \hat{u}^{(k)} $ \;
}

Decode the final state $\hat{y} \leftarrow \mathcal{D}(y^{(K)}, x)$ \;

\Return{$\hat{y}$}\;

\end{algorithm}

\paragraph{From data generation to solution generation.}


\subsection{Neural Networks and Graph Neural Networks}
\label{sec:neural_networks}

Neural networks (NNs) are a class of parametric models designed to approximate complex nonlinear mappings between inputs and outputs. They are inspired by the structure of biological neurons but are best viewed as powerful function approximators that can represent a broad class of continuous functions \cite{hornik_multilayer_1989,goodfellow_deep_2016}. A neural network defines a function 
\begin{equation}
    f_{\theta} : \mathcal{X} \rightarrow \mathcal{Y},
\end{equation}
composed of a sequence of linear transformations followed by nonlinear activation functions:
\begin{equation}
    h^{(l)} = \sigma\!\left(W^{(l)} h^{(l-1)} + b^{(l)}\right), \quad l = 1, \dots, L,
\end{equation}
where $h^{(0)} = x \in \mathcal{X}$ is the input, $h^{(L)}$ represents the output layer, and the set of trainable parameters is $\theta = \{W^{(l)}, b^{(l)}\}_{l=1}^{L}$. The activation function $\sigma$ introduces nonlinearity (e.g., ReLU, sigmoid, or tanh), which allows the network to model complex relationships that cannot be captured by purely linear models. 

During training, the parameters $\theta$ are optimized to minimize a task-specific loss function (see Equation~\ref{equ:training_task}), typically using stochastic gradient-based methods such as backpropagation and variants of gradient descent \cite{rumelhart_learning_1986,kingma_adam_2014}. The success of neural networks in computer vision, language modeling, and reinforcement learning has motivated their extension to domains where the data exhibit rich, relational or graph-structured dependencies.

\begin{figure}[H]
\centering

% Adjustable layer sizes
\def\nI{3}   % inputs
\def\nH{5}   % hidden layer 1
\def\nHtwo{4}% hidden layer 2
\def\nO{2}   % outputs

% Spacing
\def\xsep{2.6}  % horizontal separation
\def\ysep{0.9}  % vertical separation

\begin{tikzpicture}[
  >=stealth,
  node distance=\xsep,
  every node/.style={font=\small},
  neuron/.style={circle, draw, thick, minimum size=7.5mm, inner sep=0pt},
  bias/.style={circle, draw, thick, minimum size=6.2mm, inner sep=0pt, fill=gray!10},
  layerbox/.style={rounded corners, draw=black!30, inner sep=6pt}
]

% ---------- Helper to place a vertical stack of nodes ----------
% #1 = count, #2 = x, #3 = nameprefix
\newcommand{\stacknodes}[3]{%
  \pgfmathsetmacro{\N}{#1}
  \foreach \i in {1,...,#1} {%
    \node[neuron] (#3\i) at (#2, {-(\i-1)*\ysep + 0.5*(\N-1)*\ysep}) {};
  }%
}

% ---------- Place layers ----------
\stacknodes{\nI}{0}{I}                % Inputs
\node[bias] (b0) [above=5pt of I1] {$1$};

\stacknodes{\nH}{\xsep}{H}            % Hidden 1
\node[bias] (b1) [above=5pt of H1] {$1$};

\stacknodes{\nHtwo}{2*\xsep}{G}       % Hidden 2
\node[bias] (b2) [above=5pt of G1] {$1$};

\stacknodes{\nO}{3*\xsep}{O}          % Outputs

% ---------- Connections ----------
% Input -> H1
\foreach \i in {1,...,\nI} {
  \foreach \j in {1,...,\nH} {
    \draw[->, thick] (I\i) -- (H\j);
  }
}
% bias to H1
\foreach \j in {1,...,\nH} {
  \draw[->, thick] (b0) -- (H\j);
}

% H1 -> H2
\foreach \i in {1,...,\nH} {
  \foreach \j in {1,...,\nHtwo} {
    \draw[->, thick] (H\i) -- (G\j);
  }
}
% bias to H2
\foreach \j in {1,...,\nHtwo} {
  \draw[->, thick] (b1) -- (G\j);
}

% H2 -> Output
\foreach \i in {1,...,\nHtwo} {
  \foreach \j in {1,...,\nO} {
    \draw[->, thick] (G\i) -- (O\j);
  }
}
% bias to Output
\foreach \j in {1,...,\nO} {
  \draw[->, thick] (b2) -- (O\j);
}

% ---------- Labels ----------
\node[above=1.0em of b0] {\textbf{Input layer} $\;x$};
\node at ($(H3)+(0,1.35)$) {\textbf{Hidden 1} $\;\sigma\!\big(W^{(1)}x+b^{(1)}\big)$};
\node at ($(G3)+(0,1.35)$) {\textbf{Hidden 2} $\;\sigma\!\big(W^{(2)}h^{(1)}+b^{(2)}\big)$};
\node at ($(O1)+(0,1.35)$) {\textbf{Output} $\;\hat{y}=f_{\theta}(x)$};

% Optional: annotate weights
\node at ($(\xsep/2, -2.3)$) {$W^{(1)}$};
\node at ($(3*\xsep/2, -2.3)$) {$W^{(2)}$};
\node at ($(5*\xsep/2, -2.3)$) {$W^{(3)}$};

% ---------- Layer boxes (for clarity) ----------
\node[layerbox, fit=(I1) (I\nI) (b0)] {};
\node[layerbox, fit=(H1) (H\nH) (b1)] {};
\node[layerbox, fit=(G1) (G\nHtwo) (b2)] {};
\node[layerbox, fit=(O1) (O\nO)] {};

\end{tikzpicture}

\caption{Feed-forward neural network with two hidden layers, bias units, and dense connections. Activations $\sigma(\cdot)$ are applied element-wise; $W^{(l)}$ and $b^{(l)}$ denote weights and biases.}
\label{fig:nn_tikz}
\end{figure}

\paragraph{Graph Neural Networks.}
Many real-world problems can be naturally represented as graphs $\mathcal{G} = (\mathcal{V}, \mathcal{E})$, where $\mathcal{V}$ is the set of nodes and $\mathcal{E}$ the set of edges, encoding pairwise relationships or constraints. Traditional neural architectures (e.g., CNNs, RNNs) assume Euclidean input domains and therefore fail to directly exploit the topological structure of such data. \emph{Graph Neural Networks} (GNNs) address this limitation by defining message-passing operations that iteratively exchange information between neighboring nodes \cite{scarselli_graph_2009,kipf_semi-supervised_2017,bronstein_geometric_2017}. 

Each node $i \in \mathcal{V}$ is associated with an initial feature vector $h_i^{(0)} \in \mathbb{R}^{d_0}$, and the node representations are updated through successive \emph{message passing} layers of the general form:
\begin{equation}
    h_i^{(l)} = \phi^{(l)} \!\left( h_i^{(l-1)}, \, \bigoplus_{j \in \mathcal{N}(i)} \psi^{(l)}(h_i^{(l-1)}, h_j^{(l-1)}, e_{ij}) \right),
    \label{eq:gnn_message_passing}
\end{equation}
where $\mathcal{N}(i)$ denotes the neighborhood of node $i$, $e_{ij}$ represents optional edge features, $\bigoplus$ is a permutation-invariant aggregation operator (e.g., sum, mean, or max), and $\psi^{(l)}$ and $\phi^{(l)}$ are neural functions parameterized by trainable weights. This formulation ensures that GNNs are invariant to node permutations and can generalize across graphs of varying size and structure.

After several message-passing iterations, node representations can be used to produce node-level, edge-level, or global graph-level outputs, depending on the downstream task:
\begin{equation}
    y = \rho\!\left( \{ h_i^{(L)} \}_{i \in \mathcal{V}} \right),
\end{equation}
where $\rho$ is a readout function (e.g., summation or attention pooling) producing a fixed-size representation for the graph.

GNNs have emerged as a natural choice for learning in domains such as chemistry, social network analysis, combinatorial optimization, and structured prediction \cite{gilmer_neural_2017,battaglia_relational_2018}. In the context of combinatorial optimization, they offer a way to learn heuristics or inference rules that operate directly on graph-structured problem instances (e.g., traveling salesman, maximum independent set, or set covering), bridging the gap between neural models and discrete optimization techniques.

%\begin{figure}[!t]
\centering
\begin{tikzpicture}[
  >=Stealth,
  node distance=1.8cm,
  every node/.style={font=\small},
  vtx/.style={circle, draw, thick, minimum size=7.5mm, inner sep=0pt, fill=white},
  agg/.style={diamond, draw, thick, inner sep=2pt, aspect=2, fill=white},
  box/.style={rounded corners, draw=black!30, inner sep=6pt},
  msg/.style={-{Stealth[length=2.2mm,width=1.5mm]}, thick},
  edge/.style={thick},
  faded/.style={black!55}
]

% ---------- Graph nodes (layer l) ----------
\node[vtx] (v1) at (0,0) {$h^{(l)}_{1}$};
\node[vtx] (v2) [above right=1.2cm and 2.2cm of v1] {$h^{(l)}_{2}$};
\node[vtx] (v3) [right=3.8cm of v1] {$h^{(l)}_{3}$};
\node[vtx] (v4) [below right=1.2cm and 2.2cm of v1] {$h^{(l)}_{4}$};
\node[vtx,fill=white] (vi) [above right=0.0cm and 1.9cm of v1] {$h^{(l)}_{i}$};

% Optional edge features label
\node[faded] at ($(vi)!0.5!(v2)+(0,0.35)$) {$e_{2i}$};
\node[faded] at ($(vi)!0.5!(v3)+(0.2,0.25)$) {$e_{3i}$};
\node[faded] at ($(vi)!0.5!(v4)+(0,-0.35)$) {$e_{4i}$};
\node[faded] at ($(vi)!0.5!(v1)+(-0.2,0)$) {$e_{1i}$};

% ---------- Undirected edges ----------
\draw[edge] (v1) -- (vi);
\draw[edge] (v2) -- (vi);
\draw[edge] (v3) -- (vi);
\draw[edge] (v4) -- (vi);
\draw[edge,faded] (v1) -- (v2);
\draw[edge,faded] (v2) -- (v3);
\draw[edge,faded] (v3) -- (v4);
\draw[edge,faded] (v4) -- (v1);

% ---------- Message arrows into i ----------
\draw[msg] (v1) .. controls ($(v1)!0.5!(vi)+(-0.3,0.4)$) .. ($(vi)+(-0.3,0.0)$);
\draw[msg] (v2) .. controls ($(v2)!0.5!(vi)+(0.15,0.35)$) .. ($(vi)+(0.05,0.25)$);
\draw[msg] (v3) .. controls ($(v3)!0.5!(vi)+(0.25,0.0)$) .. ($(vi)+(0.35,0.0)$);
\draw[msg] (v4) .. controls ($(v4)!0.5!(vi)+(0.10,-0.35)$) .. ($(vi)+(0.05,-0.25)$);

% ---------- Aggregation operator and update ----------
\node[agg, minimum width=7mm, inner sep=1pt] (op) at ($(vi)+(1.6,0)$) {$\bigoplus$};
\draw[msg] (vi) -- (op);

% New representation at layer l+1
\node[vtx, fill=white] (vi2) [right=1.6cm of op] {$h^{(l+1)}_{i}$};
\draw[msg] (op) -- (vi2);

% ---------- Layer annotation (equations) ----------
\node[align=left, anchor=west, text width=6.5cm] (eqbox) 
  at ($(vi2)+(2.4,0)$) {$\displaystyle
  \begin{aligned}
  m^{(l)}_{j\rightarrow i} &= \psi^{(l)}\!\big(h^{(l)}_{i},\,h^{(l)}_{j},\,e_{ji}\big) \\
  \tilde{h}^{(l)}_{i} &= \bigoplus_{j\in\mathcal{N}(i)} m^{(l)}_{j\rightarrow i} \\
  h^{(l+1)}_{i} &= \phi^{(l)}\!\big(h^{(l)}_{i},\,\tilde{h}^{(l)}_{i}\big)
  \end{aligned}$};

% ---------- Readout (graph-level) ----------
% Pool nodes -> graph embedding -> output
\node[agg, minimum width=8mm] (pool) at ($(v3)+(0,-2.1)$) {$\rho$};
\draw[msg] (v1) -- (pool);
\draw[msg] (v2) -- (pool);
\draw[msg] (v3) -- (pool);
\draw[msg] (v4) -- (pool);
\draw[msg] (vi) -- (pool);

\node[vtx, minimum size=7mm] (gemb) [right=1.7cm of pool] {$\;g\;$};
\draw[msg] (pool) -- (gemb);

\node[rectangle, draw, thick, minimum width=10mm, minimum height=7mm, align=center] (out) 
  [right=1.6cm of gemb] {$y$};
\draw[msg] (gemb) -- (out);

% ---------- Labels ----------
\node[above=6pt of v2] {\textbf{Message-passing (layer $l$)}};
\node[below=4pt of pool] {\textbf{Readout}};

% ---------- Light boxes around components ----------
\node[box, fit=(v1) (v2) (v3) (v4) (vi)] (graphbox) {};
\node[box, fit=(pool) (gemb) (out)] (readbox) {};

% ---------- Legend ----------
\node[draw, rounded corners, align=left, anchor=north west, inner sep=5pt] at ($(graphbox.north west)+(-0.1,1.0)$) {%
\begin{tabular}{@{}l@{ }l@{}}
\raisebox{0.5ex}{\tikz{\node[vtx, minimum size=3.5mm]{};}} & node feature $h^{(l)}_i$ \\
\raisebox{0.5ex}{\tikz{\draw[edge] (0,0) -- (0.5,0);}} & edge with feature $e_{ij}$ \\
\raisebox{0.5ex}{\tikz{\draw[msg] (0,0) -- (0.6,0);}} & message $m^{(l)}_{j\to i}$ \\
\raisebox{0.6ex}{\tikz{\node[agg, minimum width=4.5mm, inner sep=0pt] {\small$\bigoplus$};}} & permutation-invariant aggregation
\end{tabular}
};

\end{tikzpicture}
\caption{Graph Neural Network (GNN) schematic. Messages $m^{(l)}_{j\to i}=\psi^{(l)}(h^{(l)}_{i},h^{(l)}_{j},e_{ji})$ from neighbors $j\in\mathcal{N}(i)$ are aggregated with a permutation-invariant operator $\bigoplus$ to form $\tilde{h}^{(l)}_{i}$, followed by an update $h^{(l+1)}_{i}=\phi^{(l)}(h^{(l)}_{i},\tilde{h}^{(l)}_{i})$. A readout $\rho$ pools node embeddings into a graph-level representation $g$ for prediction $y$.}
\label{fig:gnn_tikz}
\end{figure}




You said:
the last image's layout is all over the place, make it much simpler. It should fit in the width of an a4paper



ChatGPT said:
got it — here’s a much simpler, compact GNN TikZ figure that fits comfortably within an A4 page width. It shows: a tiny graph, messages into a node 
��
i, aggregation 
⨁
⨁, the updated embedding 
ℎ
��
(
��
+
1
)
h 
i
(l+1)
​
 , and a small graph-level readout. Only common TikZ libs; no big equation block.

latex
Copy code
%=== In your preamble ===
\usepackage{tikz}
\usetikzlibrary{arrows.meta,positioning,calc}

%=== Wherever you want the figure ===
\begin{figure}[!t]
\centering
\begin{tikzpicture}[
  >=Stealth,
  every node/.style={font=\small},
  vtx/.style={circle, draw, thick, minimum size=6.5mm, inner sep=0pt, fill=white},
  agg/.style={diamond, draw, thick, inner sep=1.3pt, aspect=2, fill=white},
  box/.style={rounded corners, draw=black!30, inner sep=4pt},
  msg/.style={-{Stealth[length=2mm,width=1.2mm]}, thick},
  edge/.style={thick, black!65}
]

% --- Small input graph (layer l) ---
\node[vtx] (v1) at (0,0) {$h^{(l)}_{1}$};
\node[vtx] (v2) at (1.3,1.0) {$h^{(l)}_{2}$};
\node[vtx,fill=white] (vi) at (2.6,0.0) {$h^{(l)}_{i}$};
\node[vtx] (v3) at (1.3,-1.0) {$h^{(l)}_{3}$};
\node[vtx] (v4) at (-0.2,-0.9) {$h^{(l)}_{4}$};

% Undirected edges (light)
\draw[edge] (v1) -- (v2);
\draw[edge] (v1) -- (v3);
\draw[edge] (v2) -- (vi);
\draw[edge] (v3) -- (vi);
\draw[edge] (v4) -- (v1);

% Messages into node i (simple curved arrows)
\draw[msg] (v1) .. controls ($(v1)!0.5!(vi)+(0,0.35)$) .. ($(vi)+(-0.15,0.25)$);
\draw[msg] (v2) .. controls ($(v2)!0.5!(vi)+(0.15,0.25)$) .. ($(vi)+(0.15,0.18)$);
\draw[msg] (v3) .. controls ($(v3)!0.5!(vi)+(0.10,-0.30)$) .. ($(vi)+(0.10,-0.20)$);
\draw[msg] (v4) .. controls ($(v4)!0.5!(vi)+(-0.05,-0.15)$) .. ($(vi)+(-0.20,-0.18)$);

% Aggregation operator and update
\node[agg, minimum width=6mm] (op) at ($(vi)+(1.4,0)$) {$\bigoplus$};
\draw[msg] (vi) -- (op);

\node[vtx] (vi2) at ($(op)+(1.6,0)$) {$h^{(l+1)}_{i}$};
\draw[msg] (op) -- (vi2);

% Tiny readout (graph-level)
\node[rectangle, draw, thick, minimum width=7mm, minimum height=6mm, right=10mm of vi2] (out) {$y$};
\node[draw, circle, minimum size=6.2mm, inner sep=0pt, left=6mm of out] (gemb) {$g$};
\draw[msg] (vi2) -- (gemb);
\draw[msg] (gemb) -- (out);

% Short labels (compact)
\node[above=2pt of vi] {Message passing};
\node[below=0pt of gemb] {Readout};

% Minimal legend (very compact)
\node[align=left, anchor=north west, draw, rounded corners, inner sep=3pt]
  at ([xshift=-2pt,yshift=6pt]v1.north west) {%
  \begin{tabular}{@{}l@{}}
  \tikz{\node[vtx, minimum size=4mm]{};} node $h^{(l)}$ \\
  \tikz{\draw[msg] (0,0)--(0.55,0);} message \\
  \tikz{\node[agg, minimum width=4.2mm]{\tiny$\bigoplus$};} aggregation \\
  \end{tabular}
};

\end{tikzpicture}
\caption{Compact GNN schematic. Neighbors send messages to node $i$, messages are aggregated with a permutation-invariant operator $\bigoplus$, then $h^{(l+1)}_{i}$ is produced. A simple readout pools to graph embedding $g$ and prediction $y$.}
\label{fig:gnn_tikz}
\end{figure}